% This is the template for submission of abstracts to NetSci 2026 in Boston, USA.
% It is modified from NetSci2024, which was inspired by that of NetSci2016 and that of STATPHSY25.
% The editor of the booklet reserves the right to modify your submission.

% To process this file run LaTeX2e

% ******** DO NOT EDIT ****************
\documentclass[10pt]{article}
% \usepackage{mathptmx}
\usepackage[letterpaper,margin=20mm]{geometry}
\usepackage{graphicx}
\usepackage{lmodern}
\renewcommand{\familydefault}{\sfdefault}
\pagestyle{empty}
\setlength{\parskip}{0.25\baselineskip}
\renewcommand{\title}[1]{{\noindent\large\bfseries#1\medskip\\}}
\renewcommand{\author}[2]{{\noindent #1 \medskip\\ \noindent \small #2 \medskip\\}}
% *************************************

\begin{document}

% ********** USER DEFINED *************


% Enter title here
\title{Community Detection in Labeled Growing Hypergraphs}
\author{
% Enter author(s) here
Francis Cataldo,\textsuperscript{1}
Philip Samuel Chodrow\textsuperscript{1}
}
{
% Enter affiliation(s) here
1. Middlebury College, Middlebury, VT, USA
}

We consider a model of growing hypergraphs in which in which hyperedges form as noisy copies of prior hyperedges.
Our model extends prior work [1] by incorporating binary community node labels. 
The formation of new edges is affected by homophily with respect to these labels.
Unlike standard community detection methods which assume that edges form independently conditioned on node labels, our model instead explicitly assumes that new interactions are most likely to arrive between similar nodes which have previously iteracted. 
The resulting model ``prefers'' very different community structures than those generated by stochastic blockmodel variants. 
To perform community detection under this model, we develop a likelihood-based inference method.
We show that under a continuous relaxation of the label vector, a first-order expansion of the log-likelihood function simplifies to a tractable structure which enables scalable inference via gradient methods. 
The latent-variable structure of this model also allows us to estimate the genealogy of edges, giving an approximate branching process describing the emergence of new interactions represented in the data. 

On synthetic data, we find that our method out performs standard graph-based community detection algorithms applied to the clique-projection of the hypergraph. 
We also apply our methods to two classes of node-labeled data sets which do not obey standard community detection assumptions: legislate cosponsorship in the US congress (tagged by party affiliation) and co-authorship data from arXiv (tagged by inferred author gender).

% Enter abstract here. Please keep everything within one page.

\bigskip
{\small
\noindent[1] He, Chodrow, and Mucha. (2025) ``Edge Correlations and Link Prediction in Growing Hypergraphs.'' \emph{Physical Review E}.
}

\begin{figure}[!h]
  \centering
%   \includegraphics[width=0.9\textwidth]{netsci2026.png}
  \caption{\textbf{Abstract must include a figure.}}
\end{figure}

% *************************************
\end{document}